% =============================================================================
% Appendix C: Reproducibility Protocol
% =============================================================================
\section{Reproducibility Protocol}\label{app:reproducibility}

This appendix documents the complete procedure for reproducing all
results presented in this paper.

\subsection{Software requirements}

\begin{itemize}[leftmargin=2em]
  \item \textbf{Python} $\geq$ 3.10 (tested with 3.14.0).
  \item \textbf{NumPy} $\geq$ 2.0 (tested with 2.2.6).
    NumPy is the only external dependency.
  \item \textbf{LaTeX} distribution with pdfLaTeX (tested with
    TeX~Live 2024), including packages: \texttt{pgfplots} ($\geq$
    1.18), \texttt{tikz}, \texttt{siunitx}, \texttt{longtable},
    \texttt{booktabs}, \texttt{natbib}, \texttt{cleveref},
    \texttt{threeparttable}, \texttt{pdflscape}, \texttt{adjustbox},
    \texttt{subcaption}, \texttt{listings}, \texttt{xcolor}.
\end{itemize}

\subsection{Directory structure}

The project is organised as follows:

\begin{verbatim}
paper2/
+-- main.tex                 % Master document
+-- references.bib           % Bibliography
+-- sections/
|   +-- 00_executive_summary.tex
|   +-- 01_introduction.tex
|   +-- ...
|   +-- 10_data_availability.tex
+-- tables/                  % Generated .tex table fragments
|   +-- T06_axis_top_bottom.tex
|   +-- T07_axis_dominance.tex
|   +-- ...
|   +-- radar_*.dat
+-- figures/                 % TikZ figure source files
|   +-- F01_histogram.tex
|   +-- ...
|   +-- F15_radar_FR.tex
+-- appendix/
|   +-- appendix_a_tables.tex
|   +-- appendix_b_robustness.tex
|   +-- appendix_c_reproducibility.tex
|   +-- appendix_d_code.tex
+-- scripts/
    +-- isi_results_paper2.py
\end{verbatim}

\subsection{Step-by-step reproduction}

\begin{enumerate}[leftmargin=2em]
  \item \textbf{Install Python and NumPy.}
    \begin{verbatim}
    python -m venv .venv
    source .venv/bin/activate
    pip install numpy
    \end{verbatim}

  \item \textbf{Run the computation script.}
    \begin{verbatim}
    cd paper2/scripts
    python isi_results_paper2.py
    \end{verbatim}
    This generates all \texttt{.tex} table fragments and
    \texttt{.dat} coordinate files in the \texttt{tables/} directory.
    The script embeds the full ISI~v1.0 data snapshot and performs
    all computations (LOO, z-score, Dirichlet MC, winsorization,
    variance decomposition, Lorenz curve, ECDF, histograms, radar
    coordinates).

  \item \textbf{Compile the LaTeX document.}
    \begin{verbatim}
    cd paper2
    pdflatex main.tex
    bibtex main
    pdflatex main.tex
    pdflatex main.tex
    \end{verbatim}
    Three passes are required to resolve cross-references and the
    bibliography.

  \item \textbf{Verify.}  The final PDF should compile with
    0~errors, 0~warnings (excluding font-related warnings), and
    0~overfull boxes.
\end{enumerate}

\subsection{Random seed}

The Dirichlet Monte Carlo exercise uses a fixed random seed
(\texttt{np.random.seed(42)}) to ensure exact reproducibility of
rank volatility statistics.  Changing the seed will produce
quantitatively similar but not identical results.

\subsection{Data integrity}

The ISI~v1.0 data snapshot is embedded directly in the Python script
(\texttt{isi\_results\_\allowbreak paper2.py}, lines 37--77) as a list of tuples.
No external data files need to be downloaded or configured.  The
embedded data constitutes the single authoritative source for all
computations in this paper.
